%%%%%%%%%%%%%%%%%%%%%%%%%%%%%%%%%%%%%%%%%%%%%%%%%%%%%%%%%%%%%%%%%%%%%%%%%%%%%%%%%%%%
% PREÁMBULO MAESTRO PARA "EL CÓDIGO DE LA REALIDAD"
% WASAFF CONSULTING
%%%%%%%%%%%%%%%%%%%%%%%%%%%%%%%%%%%%%%%%%%%%%%%%%%%%%%%%%%%%%%%%%%%%%%%%%%%%%%%%%%%%

% --- CONFIGURACIONES ESENCIALES ---
\usepackage[utf8]{inputenc}   % Acentos y caracteres en español
\usepackage[T1]{fontenc}      % Codificación de fuentes moderna
\usepackage[spanish,es-noshorthands]{babel} % Soporte para español
\usepackage{lmodern}          % Fuente Latin Modern (mejora de Computer Modern)

% --- GEOMETRÍA Y DISEÑO DE PÁGINA ---
% Márgenes definidos para una buena legibilidad en A4
\usepackage[a4paper, top=2.5cm, bottom=3cm, left=2.5cm, right=2.5cm]{geometry}

% --- PAQUETES MATEMÁTICOS ---
\usepackage{amsmath}          % El paquete principal de la AMS para matemáticas
\usepackage{amsfonts}         % Fuentes matemáticas
\usepackage{amssymb}          % Símbolos matemáticos
\usepackage{mathtools}        % Extensiones y correcciones para amsmath
\usepackage{bm}               % Para poner símbolos matemáticos en negrita (vectores, etc.)

% --- TIPOGRAFÍA Y ESTILO ---
\usepackage{microtype}        % Mejora la justificación y el espaciado entre caracteres
\usepackage{graphicx}         % Para incluir imágenes
\usepackage{booktabs}         % Para tablas de calidad profesional
\usepackage{xcolor}           % Para definir y usar colores
\usepackage{lipsum}           % Para generar texto de relleno (lorem ipsum)
\usepackage{csquotes}         % Para comillas inteligentes y compatibilidad con biblatex

% Paquetes adicionales necesarios
\usepackage{tcolorbox}
\tcbuselibrary{skins, breakable}
\usepackage{tikz}
\usetikzlibrary{shapes, arrows, positioning, calc}
\usepackage{pgfplots}
\pgfplotsset{compat=1.18}

% --- CÓDIGO FUENTE ---
% Minted es el estándar de oro para código con sintaxis coloreada.
% ¡IMPORTANTE! Requiere compilar con la opción -shell-escape
% Ejemplo: latexmk -pdf -shell-escape LibroMaestro.tex
\usepackage{minted}

% --- BIBLIOGRAFÍA ---
% Usamos biblatex y biber, el sistema moderno y más flexible.
\usepackage[backend=biber, style=numeric, sorting=none]{biblatex}
\addbibresource{Bibliografia/BibliografiaMaestra.bib} % Apunta a nuestro archivo .bib

% --- ENLACES Y REFERENCIAS ---
% hyperref debe ser uno de los últimos paquetes en cargarse.
\usepackage{hyperref}
\hypersetup{
    colorlinks=true,
    linkcolor=blue,        % Color para enlaces internos (índice, citas)
    citecolor=green,       % Color para citas bibliográficas
    filecolor=magenta,
    urlcolor=cyan,         % Color para URLs
    pdftitle={El Código de la Realidad},
    pdfauthor={Felipe Wasaff},
    bookmarks=true,
    pdfpagemode=FullScreen,
}

% --- INFORMACIÓN DEL TÍTULO ---
% Este bloque es el que soluciona el error "! LaTeX Error: No \title given."
\title{El Código de la Realidad \\ \large Simulación Multifísica desde los Primeros Principios para la Ingeniería Industrial}
\author{Felipe Wasaff \\ Wasaff Consulting}
\date{\today}

% --- COMANDOS PERSONALIZADOS Y NUEVOS ENTORNOS ---
\usepackage{newfloat} % Para crear nuevos tipos de flotantes
% Aquí está la corrección: cambiamos 'listing' por 'codelisting'
\DeclareFloatingEnvironment[fileext=lol, listname={Índice de Listados de Código}, name=Listado de Código, within=chapter]{codelisting}

\newcommand{\fenics}{\textsc{FEniCS}}